\documentclass[reqno]{amsart}
\usepackage[T1]{fontenc}
\usepackage{amsmath,amssymb,amsthm}
\usepackage{xr-hyper}
\usepackage[hidelinks]{hyperref}
\externaldocument[sm-]{20261125T000000--slope-method__workbook}
\externaldocument[cv-]{20261125T006000--covolumes__workbook}

% Chapter number; \ref{chap} gives it.
\newcounter{chap}
% Chapter numbers, generated from 20261125T012000--grh-workbook__guide.org
% by scripts/make-chapters.py.  Do not edit: change the order of the
% index and rerun the script (./preview-site.sh and CI do this).
%
% Usage in a sheet, after \newcounter{chap}:
%   % Chapter numbers, generated from 20261125T012000--grh-workbook__guide.org
% by scripts/make-chapters.py.  Do not edit: change the order of the
% index and rerun the script (./preview-site.sh and CI do this).
%
% Usage in a sheet, after \newcounter{chap}:
%   % Chapter numbers, generated from 20261125T012000--grh-workbook__guide.org
% by scripts/make-chapters.py.  Do not edit: change the order of the
% index and rerun the script (./preview-site.sh and CI do this).
%
% Usage in a sheet, after \newcounter{chap}:
%   \input{chapters}\chapnumber{<sheet ID>}
\makeatletter
\providecommand\chapters@set[2]{\@namedef{chapters@#1}{#2}}
\chapters@set{20261125T001000}{1}% zeros-to-ramp-sums
\chapters@set{20261125T001100}{2}% inert-primes-and-estimate-b
\chapters@set{20261125T005000}{3}% galois-fourier-and-averaging
\chapters@set{20261125T005500}{4}% columns-from-a-sphere
\chapters@set{20261125T000000}{5}% slope-method
\chapters@set{20261126T030100}{6}% a-twist-for-each-prime
\chapters@set{20261125T003000}{7}% the-chain-and-the-square-root
\chapters@set{20261125T010000}{8}% rank-lemma-and-zero-estimate
\chapters@set{20261125T010500}{9}% poincare-series-and-near-orthogonality
\chapters@set{20261124T223000}{10}% elliptic-curves-for-the-proof
\chapters@set{20261125T009500}{11}% fock-model-and-theta-intertwiners
\chapters@set{20261124T223100}{12}% complex-multiplication-and-inert-reduction
\chapters@set{20261125T008000}{13}% one-cm-determinant
\chapters@set{20261124T214500}{14}% supersingular-theta-divisibility
\chapters@set{20261124T214600}{15}% the-isotropic-plane
\chapters@set{20261125T013700}{16}% pairings-on-a-formal-group-of-height-two
\chapters@set{20261125T013500}{17}% hecke-nondegeneracy-on-the-tree
\chapters@set{20261126T030200}{18}% the-bound-is-attained
\chapters@set{20261125T009000}{19}% kronecker-limit-and-chowla-selberg
\chapters@set{20261125T002000}{20}% definite-theta-lifts
\chapters@set{20261125T015500}{21}% siegel-weil-in-the-smallest-case
\chapters@set{20261125T002100}{22}% definite-theta-lifts-and-norms
\chapters@set{20261125T016500}{23}% hecke-operators-and-theta-lifts
\chapters@set{20261125T016700}{24}% divisibility-along-the-chain
\chapters@set{20261125T004000}{25}% denominators-and-sha
\chapters@set{20261125T004100}{26}% tate-twists-and-sha
\chapters@set{20261126T020430}{27}% the-minus-part-by-gauss-sums
\chapters@set{20261125T014100}{28}% congruence-modules
\chapters@set{20261125T014500}{29}% eisenstein-congruences-and-selmer-classes
\chapters@set{20261125T014700}{30}% pseudocharacters-and-the-length-bound
\chapters@set{20261125T011000}{31}% galois-representations-black-box
\chapters@set{20261126T030300}{32}% the-data
\chapters@set{20261126T020210}{33}% bessel-poincare-series-and-their-gram-operators
\chapters@set{20261126T030400}{34}% the-seesaw-identity
\chapters@set{20261126T030500}{35}% divisibility-of-the-composite-lift
\chapters@set{20261126T030600}{36}% theorem-a-short-forms
\chapters@set{20261126T030700}{37}% the-determinant-of-a-chain
\chapters@set{20261126T030800}{38}% the-cost-of-one-lift
\chapters@set{20261126T030900}{39}% denominators-of-projections-to-lifts
\chapters@set{20261126T031000}{40}% twists-with-few-congruences
\chapters@set{20261126T031100}{41}% theorem-c-and-the-main-inequality
\providecommand\chapnumber[1]{%
  \@ifundefined{chapters@#1}%
    {\PackageError{chapters}{Sheet #1 is not listed in the index}%
      {List it in 20261125T012000--grh-workbook__guide.org\MessageBreak
       and run scripts/make-chapters.py.}}%
    {\edef\thechap{\@nameuse{chapters@#1}}}}
\makeatother
\chapnumber{<sheet ID>}
\makeatletter
\providecommand\chapters@set[2]{\@namedef{chapters@#1}{#2}}
\chapters@set{20261125T001000}{1}% zeros-to-ramp-sums
\chapters@set{20261125T001100}{2}% inert-primes-and-estimate-b
\chapters@set{20261125T005000}{3}% galois-fourier-and-averaging
\chapters@set{20261125T005500}{4}% columns-from-a-sphere
\chapters@set{20261125T000000}{5}% slope-method
\chapters@set{20261126T030100}{6}% a-twist-for-each-prime
\chapters@set{20261125T003000}{7}% the-chain-and-the-square-root
\chapters@set{20261125T010000}{8}% rank-lemma-and-zero-estimate
\chapters@set{20261125T010500}{9}% poincare-series-and-near-orthogonality
\chapters@set{20261124T223000}{10}% elliptic-curves-for-the-proof
\chapters@set{20261125T009500}{11}% fock-model-and-theta-intertwiners
\chapters@set{20261124T223100}{12}% complex-multiplication-and-inert-reduction
\chapters@set{20261125T008000}{13}% one-cm-determinant
\chapters@set{20261124T214500}{14}% supersingular-theta-divisibility
\chapters@set{20261124T214600}{15}% the-isotropic-plane
\chapters@set{20261125T013700}{16}% pairings-on-a-formal-group-of-height-two
\chapters@set{20261125T013500}{17}% hecke-nondegeneracy-on-the-tree
\chapters@set{20261126T030200}{18}% the-bound-is-attained
\chapters@set{20261125T009000}{19}% kronecker-limit-and-chowla-selberg
\chapters@set{20261125T002000}{20}% definite-theta-lifts
\chapters@set{20261125T015500}{21}% siegel-weil-in-the-smallest-case
\chapters@set{20261125T002100}{22}% definite-theta-lifts-and-norms
\chapters@set{20261125T016500}{23}% hecke-operators-and-theta-lifts
\chapters@set{20261125T016700}{24}% divisibility-along-the-chain
\chapters@set{20261125T004000}{25}% denominators-and-sha
\chapters@set{20261125T004100}{26}% tate-twists-and-sha
\chapters@set{20261126T020430}{27}% the-minus-part-by-gauss-sums
\chapters@set{20261125T014100}{28}% congruence-modules
\chapters@set{20261125T014500}{29}% eisenstein-congruences-and-selmer-classes
\chapters@set{20261125T014700}{30}% pseudocharacters-and-the-length-bound
\chapters@set{20261125T011000}{31}% galois-representations-black-box
\chapters@set{20261126T030300}{32}% the-data
\chapters@set{20261126T020210}{33}% bessel-poincare-series-and-their-gram-operators
\chapters@set{20261126T030400}{34}% the-seesaw-identity
\chapters@set{20261126T030500}{35}% divisibility-of-the-composite-lift
\chapters@set{20261126T030600}{36}% theorem-a-short-forms
\chapters@set{20261126T030700}{37}% the-determinant-of-a-chain
\chapters@set{20261126T030800}{38}% the-cost-of-one-lift
\chapters@set{20261126T030900}{39}% denominators-of-projections-to-lifts
\chapters@set{20261126T031000}{40}% twists-with-few-congruences
\chapters@set{20261126T031100}{41}% theorem-c-and-the-main-inequality
\providecommand\chapnumber[1]{%
  \@ifundefined{chapters@#1}%
    {\PackageError{chapters}{Sheet #1 is not listed in the index}%
      {List it in 20261125T012000--grh-workbook__guide.org\MessageBreak
       and run scripts/make-chapters.py.}}%
    {\edef\thechap{\@nameuse{chapters@#1}}}}
\makeatother
\chapnumber{<sheet ID>}
\makeatletter
\providecommand\chapters@set[2]{\@namedef{chapters@#1}{#2}}
\chapters@set{20261125T001000}{1}% zeros-to-ramp-sums
\chapters@set{20261125T001100}{2}% inert-primes-and-estimate-b
\chapters@set{20261125T005000}{3}% galois-fourier-and-averaging
\chapters@set{20261125T005500}{4}% columns-from-a-sphere
\chapters@set{20261125T000000}{5}% slope-method
\chapters@set{20261126T030100}{6}% a-twist-for-each-prime
\chapters@set{20261125T003000}{7}% the-chain-and-the-square-root
\chapters@set{20261125T010000}{8}% rank-lemma-and-zero-estimate
\chapters@set{20261125T010500}{9}% poincare-series-and-near-orthogonality
\chapters@set{20261124T223000}{10}% elliptic-curves-for-the-proof
\chapters@set{20261125T009500}{11}% fock-model-and-theta-intertwiners
\chapters@set{20261124T223100}{12}% complex-multiplication-and-inert-reduction
\chapters@set{20261125T008000}{13}% one-cm-determinant
\chapters@set{20261124T214500}{14}% supersingular-theta-divisibility
\chapters@set{20261124T214600}{15}% the-isotropic-plane
\chapters@set{20261125T013700}{16}% pairings-on-a-formal-group-of-height-two
\chapters@set{20261125T013500}{17}% hecke-nondegeneracy-on-the-tree
\chapters@set{20261126T030200}{18}% the-bound-is-attained
\chapters@set{20261125T009000}{19}% kronecker-limit-and-chowla-selberg
\chapters@set{20261125T002000}{20}% definite-theta-lifts
\chapters@set{20261125T015500}{21}% siegel-weil-in-the-smallest-case
\chapters@set{20261125T002100}{22}% definite-theta-lifts-and-norms
\chapters@set{20261125T016500}{23}% hecke-operators-and-theta-lifts
\chapters@set{20261125T016700}{24}% divisibility-along-the-chain
\chapters@set{20261125T004000}{25}% denominators-and-sha
\chapters@set{20261125T004100}{26}% tate-twists-and-sha
\chapters@set{20261126T020430}{27}% the-minus-part-by-gauss-sums
\chapters@set{20261125T014100}{28}% congruence-modules
\chapters@set{20261125T014500}{29}% eisenstein-congruences-and-selmer-classes
\chapters@set{20261125T014700}{30}% pseudocharacters-and-the-length-bound
\chapters@set{20261125T011000}{31}% galois-representations-black-box
\chapters@set{20261126T030300}{32}% the-data
\chapters@set{20261126T020210}{33}% bessel-poincare-series-and-their-gram-operators
\chapters@set{20261126T030400}{34}% the-seesaw-identity
\chapters@set{20261126T030500}{35}% divisibility-of-the-composite-lift
\chapters@set{20261126T030600}{36}% theorem-a-short-forms
\chapters@set{20261126T030700}{37}% the-determinant-of-a-chain
\chapters@set{20261126T030800}{38}% the-cost-of-one-lift
\chapters@set{20261126T030900}{39}% denominators-of-projections-to-lifts
\chapters@set{20261126T031000}{40}% twists-with-few-congruences
\chapters@set{20261126T031100}{41}% theorem-c-and-the-main-inequality
\providecommand\chapnumber[1]{%
  \@ifundefined{chapters@#1}%
    {\PackageError{chapters}{Sheet #1 is not listed in the index}%
      {List it in 20261125T012000--grh-workbook__guide.org\MessageBreak
       and run scripts/make-chapters.py.}}%
    {\edef\thechap{\@nameuse{chapters@#1}}}}
\makeatother
\chapnumber{20261125T017500}
\renewcommand{\thesection}{\thechap.\arabic{section}}
\renewcommand{\theequation}{\thechap.\arabic{equation}}

\theoremstyle{definition}
\newtheorem{problem}{Problem}[section]
\theoremstyle{plain}
\newtheorem{theorem}{Theorem}[section]

\newcommand{\Z}{\mathbb{Z}}
\newcommand{\Q}{\mathbb{Q}}
\newcommand{\R}{\mathbb{R}}
\newcommand{\C}{\mathbb{C}}
\newcommand{\Pb}{\mathbb{P}}
\newcommand{\ii}{\mathrm{i}}
\newcommand{\covol}{\operatorname{covol}}
\newcommand{\disc}{\operatorname{disc}}
\newcommand{\mS}{m_{S}}
\newcommand{\hFS}{h_{\mathrm{FS}}}

\setlength{\emergencystretch}{2em}

\title{Chow forms and an effective Hilbert--Samuel bound}
\author{}
\date{November 25, 2026}

\begin{document}
\maketitle
\refstepcounter{chap}\label{chap}

\paragraph{Prerequisites.}
Jensen's formula; Gauss's lemma for polynomials over a factorial ring; binary forms and their discriminants; B\'ezout's theorem for plane curves.  Mahler measures \cite[\S1.6]{BG} and Chow forms \cite[Chs.~3--4]{GKZ} are recalled in the text.  Hadamard's inequality \(\lvert\det(v_1, \dots, v_r)\rvert \le \prod_j \lVert v_j \rVert\) for vectors of a Hermitian space of dimension \(r\) (for integer matrices, Problem~\ref{sm-pr:hadamard-det}).  From Chapter~\ref{sm-chap}, \emph{The slope method}: Section~\ref{sm-sec:slopes} (degrees; over \(\Q\), \(\deg L = -\log\covol L\)).  From Chapter~\ref{cv-chap}, \emph{Covolumes of lattices of sections}: the lattice \(E_D\) and Problems~\ref{cv-pr:degED} and~\ref{cv-pr:bergman}, and Theorem~\ref{cv-thm:P1}, of Section~\ref{cv-sec:P1}, and Theorem~\ref{cv-thm:HS} and formula~\eqref{cv-eq:hPN} of Section~\ref{cv-sec:HS}.

In Chapter~\ref{cv-chap} the covolume of \(E_D = H^0(\Pb^1_\Z, \mathcal O(D))\) was computed exactly because the monomials are orthogonal.  On a curve \(C \subset \Pb^2\) no orthogonal basis of sections is known, and the arithmetic Hilbert--Samuel theorem gives \(\deg H^0(C, \mathcal O(D))\) only up to \(o(D^2)\), with an error that is not explicit.  How can one bound this degree from below with every constant explicit?  The reconstruction \cite[proof dependency map, item HSAt]{Adr} describes the covolume bound of the argument as a Chow-form determinant argument in the style of Nesterenko and Philippon \cite{NP}, with Mahler heights, Hadamard's inequality and a monomial count.  This chapter treats the ingredients: Mahler measures, Chow forms as heights, the exact covolume for points on \(\Pb^1\), and a determinant of sections of \(\Pb^1\), where one sees what Hadamard's inequality loses.

\section{Mahler measures}\label{sec:mahler}

For a nonzero \(P \in \C[x_1, \dots, x_N]\) let
\[
  m(P) := \int_{[0,1]^N} \log\bigl\lvert P(e^{2\pi\ii\theta_1}, \dots, e^{2\pi\ii\theta_N})\bigr\rvert \, d\theta , \qquad M(P) := e^{m(P)} ,
\]
the (logarithmic) Mahler measure \cite[\S1.6]{BG}.  For \(G\) homogeneous on \(\C^N\), or multihomogeneous on \(\C^{N_1} \times \dots \times \C^{N_k}\), let \(\mS(G)\) be the mean of \(\log\lvert G \rvert\) over the product of the unit spheres \(S^{2N_j - 1}\), each with its invariant probability measure.  Write \(\lvert\cdot\rvert_2\) for the Euclidean norm and \(H_k := \sum_{j=1}^{k} 1/j\), with \(H_0 = 0\).

\begin{problem}\label{pr:jensen}
\(M(a \prod_i (x - \alpha_i)) = \lvert a \rvert \prod_i \max(1, \lvert\alpha_i\rvert)\), and \(M(PQ) = M(P)M(Q)\).  Hence if a binary form factors as \(F = \prod_i (\beta_i x - \alpha_i y)\), then \(M(F(x,1)) = \prod_i \max(\lvert\alpha_i\rvert, \lvert\beta_i\rvert)\).
\end{problem}

\begin{problem}\label{pr:lex}
Order the monomials of \(P \ne 0\) lexicographically with \(x_N\) most significant: compare the exponents of \(x_N\) first, then those of \(x_{N-1}\), and so on.  Let \(a\) be the coefficient of the largest monomial of \(P\).  Then \(m(P) \ge \log\lvert a \rvert\).  In particular \(M(P) \ge 1\) for every nonzero \(P \in \Z[x_1, \dots, x_N]\).
\end{problem}

\begin{problem}\label{pr:beta}
Let \(z\) be uniform on the unit sphere of \(\C^N\), \(N \ge 2\).  Then \(\lvert z_1 \rvert^2\) has density \((N-1)(1-t)^{N-2}\) on \([0,1]\), and \(\mathbb E \log\lvert z_1 \rvert^2 = -H_{N-1}\).
\end{problem}

The measure \(\mS\) on spheres was introduced by Philippon \cite{Ph1}.  Lelong \cite{Le} proved that \(0 \le m(G) - \mS(G) \le \frac12 H_{N-1}\, k\) for every \(G \in \C[u_1, \dots, u_N]\) homogeneous of degree \(k\); see \cite[(1.2)]{KPS}.  We need the lower bound in the following form.

\begin{problem}\label{pr:sphere-torus}
If \(G \in \Z[u_1, \dots, u_N]\) is nonzero and homogeneous of degree \(k\), then \(\mS(G) \ge -\tfrac12 H_{N-1}\, k\).
\end{problem}

\begin{problem}\label{pr:sphere-linear}
For \(p \in \C^N \setminus 0\), \(\mS(u_1p_1 + \dots + u_Np_N) = \log\lvert p \rvert_2 - \tfrac12 H_{N-1}\).
\end{problem}

\section{Points on the projective line}\label{sec:points}

The smallest case is a hypersurface in \(\Pb^1\), a finite set of points, where the covolume is exact.  Let \(F \in \Z[x, y]\) be a primitive binary form of degree \(\delta \ge 1\) with \(\disc F \ne 0\), and factor \(F = \prod_{i=1}^{\delta} (\omega_{i1}x - \omega_{i0}y)\) with \(\omega_i \in \C^2\).  The \(\omega_i\) are determined up to \(\omega_i \mapsto t_i\omega_i\) with \(\prod t_i = 1\), and the points \(X = \{(\omega_{i0} : \omega_{i1})\}\) are the zeros of \(F\).  The polynomial
\[
  f_X(u_0, u_1) := \prod_i (u_0\omega_{i0} + u_1\omega_{i1}) = F(u_1, -u_0)
\]
vanishes exactly when the hyperplane \(u_0x + u_1y = 0\) meets \(X\); it is the \emph{Chow form} of \(X\).

\begin{problem}\label{pr:point}
We have
\[
  \mS(f_X) = \sum_i \log\lvert\omega_i\rvert_2 - \tfrac12\delta , \qquad m(F) = \sum_i \log\max(\lvert\omega_{i0}\rvert, \lvert\omega_{i1}\rvert) ,
\]
so that \(m(F) \le \mS(f_X) + \tfrac12\delta \le m(F) + \tfrac12\delta\log 2\).  For \(\delta = 1\) and \(F = qx - py\), \(m(F)\) is the Weil height of \((p : q)\).
\end{problem}

For \(D \ge \delta - 1\) let \(\Lambda_D := \Z[x,y]_D / F\,\Z[x,y]_{D-\delta}\), the lattice of sections of \(\mathcal O(D)\) on \(X\), with
\[
  \lVert g \rVert^2 := \sum_{i=1}^{\delta} \frac{\lvert g(\omega_i) \rvert^2}{\lvert\omega_i\rvert_2^{2D}} .
\]

\begin{problem}\label{pr:quotient0}
\(F\,\Z[x,y]_{D-\delta} = \Z[x,y]_D \cap F\,\Q[x,y]_{D-\delta}\), so \(\Lambda_D\) is free of rank \(\delta\), and \(\lVert\cdot\rVert\) is a Euclidean norm on \(\Lambda_D \otimes \R\).  For linearly independent \(g_1, \dots, g_\delta \in \Lambda_D\),
\[
  \lvert\det(g_j(\omega_i))\rvert \big/ \textstyle\prod_i \lvert\omega_i\rvert_2^{D}
\]
is the covolume of the lattice they span.
\end{problem}

Let \(a\) be the coefficient of \(x^\delta\) in \(F\), and suppose \(a \ne 0\).  Let \(V := \prod_{i < i'} (\omega_{i'0}\omega_{i1} - \omega_{i0}\omega_{i'1})\).

\begin{problem}\label{pr:index0}
The images of \(x^jy^{D-j}\), \(0 \le j < \delta\), span a sublattice of index \(\lvert a \rvert^{D-\delta+1}\) in \(\Lambda_D\).
\end{problem}

\begin{problem}\label{pr:vdm}
\(\det\bigl(\omega_{i0}^{j}\omega_{i1}^{D-j}\bigr)_{1 \le i \le \delta,\, 0 \le j < \delta} = \pm a^{D-\delta+1} V\), and \(V^2 = \disc F\), where \(\disc F = a^{2\delta-2}\prod_{i<i'}(\alpha_i - \alpha_{i'})^2\) for \(F = a\prod_i (x - \alpha_i y)\).
\end{problem}

So for a basis \(g_j\) of \(\Lambda_D\) the determinant \(\det(g_j(\omega_i))\) is \(\pm V\): the index of Problem~\ref{pr:index0} cancels the power of \(a\) in Problem~\ref{pr:vdm}.  When \(a = 0\) one first applies a substitution in \(\mathrm{GL}_2(\Z)\), which changes \(\det(g_j(\omega_i))\) and \(V\) at most by a sign.  With Problems~\ref{pr:point} and~\ref{pr:quotient0} this proves the following.

\begin{theorem}\label{thm:points}
For \(D \ge \delta - 1\),
\[
  \deg \Lambda_D = D \sum_i \log\lvert\omega_i\rvert_2 - \tfrac12\log\lvert\disc F\rvert = D\bigl(\mS(f_X) + \tfrac12\delta\bigr) - \tfrac12\log\lvert\disc F\rvert .
\]
\end{theorem}

\begin{problem}\label{pr:HS0}
For \(F = x^2 + y^2\), \(\deg\Lambda_D = (D-1)\log 2\).  For \(F = x^2 - 2y^2\), \(\deg\Lambda_D = D\log 3 - \tfrac32\log 2\); for \(D = 1\) this is \(0.0588\ldots\)
\end{problem}

The slope grows linearly in \(D\) with coefficient \(\frac1\delta\bigl(\mS(f_X) + \frac12\delta\bigr)\), a Mahler measure of the Chow form plus a constant.  For a curve there is no such closed form.

\section{Chow forms}\label{sec:chow}

Let \(X \subset \Pb^n_\Q\) be irreducible of dimension \(d\) and degree \(\delta\).  The \((d+1)\)-tuples of hyperplanes \(u^{(0)}, \dots, u^{(d)}\) with a common zero on \(X(\C)\) form a hypersurface defined over \(\Q\) and irreducible over \(\Q\).  Its equation \(f_X \in \Z[u^{(0)}, \dots, u^{(d)}]\), taken primitive, is irreducible over \(\Q\) and multihomogeneous of degree \(\delta\) in each block; it is the Chow form of \(X\), unique up to sign \cite[Ch.~3, \S2 and Ch.~4, \S1]{GKZ}.  For a hypersurface \(X = V(F)\), \(f_X\) is \(F\) evaluated at the vector of signed \(n \times n\) minors of \((u^{(0)}, \dots, u^{(n-1)})\), the common zero of the \(n\) hyperplanes.  Philippon used Mahler measures of \(f_X\) as heights \cite{Ph1}.  He also proved that the Fubini--Study height of Bost, Gillet and Soul\'e \cite{BGS}, the arithmetic self-intersection of \(\bar{\mathcal O}(1)\) on \(X\), is
\begin{equation}\label{eq:philippon}
  \hFS(X) = \mS(f_X) + \tfrac12 (d+1)\,\delta\, H_n
\end{equation}
\cite{Ph3}.  By Problem~\ref{pr:point}, Theorem~\ref{thm:points} reads \(\deg\Lambda_D = D\hFS(X) - \frac12\log\lvert\disc F\rvert\).

From now on \(C = V(F) \subset \Pb^2\), with \(F \in \Z[x_0, x_1, x_2]\) primitive and irreducible of degree \(\delta\).  Lines are \(s(x) = s_0x_0 + s_1x_1 + s_2x_2\), and \(s \times l\) is the cross product.

\begin{problem}\label{pr:chow-curve}
\(f_C(s, l) = F(s \times l)\).  If the line \(s = 0\) meets \(C\) in \(\delta\) distinct points \(P_1, \dots, P_\delta \in \C^3\), then \(f_C(s, l) = c(s)\prod_j l(P_j)\) with \(c(s) \ne 0\).
\end{problem}

The hypothesis holds for \(s\) outside a proper Zariski-closed set: through a point \(Q \notin C\) only finitely many lines are tangent to \(C\) or pass through a singular point, by B\'ezout applied to \(C\) and the polar curve \(\sum_k Q_k\,\partial F/\partial x_k = 0\).

\begin{problem}\label{pr:chow-line}
For the line \(x_2 = 0\) in \(\Pb^2\), \(f = s_0l_1 - s_1l_0\) and \(\mS(f) = -1\) on \(S^5 \times S^5\).  For \(\Pb^1\) itself, with \(f = s_0l_1 - s_1l_0\) on \(S^3 \times S^3\), \(\mS(f) = -\frac12\).  In both cases~\eqref{eq:philippon} gives \(\frac12\), the value of formula~\eqref{cv-eq:hPN} of Chapter~\ref{cv-chap} for \(N = 1\).
\end{problem}

Now let the line move, so that the model changes while the curve over \(\C\) stays a projective line.  For an integer \(N \ge 0\) let \(C_N \subset \Pb^2\) be the line \(x_2 = Nx_0\); thus \(C_0\) is the line of Problem~\ref{pr:chow-line}.  Let \(H^0(C_N, \mathcal O(T))\) be the lattice of restrictions of \(\Z[x_0, x_1, x_2]_T\) to \(C_N\), with the norm
\[
  \lVert g \rVert^2 := \int_{C_N(\C)} \frac{\lvert g(z) \rvert^2}{\lvert z \rvert_2^{2T}} \, d\mu_N(z) ,
\]
where \(\mu_N\) is the probability measure of the metric induced by the Fubini--Study metric of \(\Pb^2\).  For \(N = 0\) this is the lattice \(E_T\) of Chapter~\ref{cv-chap}.

\begin{problem}\label{pr:line-deg}
The restrictions of \(x_0^ix_1^{T-i}\), \(0 \le i \le T\), form an orthogonal \(\Z\)-basis of \(H^0(C_N, \mathcal O(T))\), and
\[
  \deg H^0(C_N, \mathcal O(T)) - \deg H^0(C_0, \mathcal O(T)) = \tfrac14 T(T+1)\log(1 + N^2) .
\]
\end{problem}

\begin{problem}\label{pr:line-chow}
\(f_{C_N}(s, l) = (s \times l)_2 - N(s \times l)_0\), which is \(\det(p, s, l)\) for \(p = (-N, 0, 1)\), and \(\hFS(C_N) - \hFS(C_0) = \frac12\log(1 + N^2)\).
\end{problem}

By Theorem~\ref{cv-thm:HS} with \(d = 1\), the coefficient of \(T^2\) in \(\deg H^0(C_N, \mathcal O(T))\) is \(\frac12\hFS(C_N)\).  The change of height in Problem~\ref{pr:line-chow} therefore predicts a change \(\frac14\log(1 + N^2)\) in that coefficient, which is the coefficient of \(T^2\) in Problem~\ref{pr:line-deg}.  The height determines the leading term exactly; the term \(\frac14 T\log(1 + N^2)\) is invisible to it.

\section{A determinant of sections of the projective line}\label{sec:check}

What does Hadamard's inequality lose when the sections are not orthogonal?  Take \(X = \Pb^1\), so that \(\Lambda_D = E_D\) with the norm of Chapter~\ref{cv-chap}.  For \(s, l \in \C^2\) write also \(s = s_0x + s_1y\) and \(l = l_0x + l_1y\), and consider the sections \(s^al^{D-a}\), \(0 \le a \le D\).

\begin{problem}\label{pr:P1-det}
In the monomial basis, \(\det\bigl(s^al^{D-a}\bigr)_{0 \le a \le D} = \pm(s_0l_1 - s_1l_0)^{D(D+1)/2}\).
\end{problem}

\begin{problem}\label{pr:P1-exact}
If \(s\) and \(l\) are orthonormal, the columns are orthogonal in \(E_D \otimes \C\), Hadamard's inequality is an equality, and it gives Problem~\ref{cv-pr:degED} again.
\end{problem}

\begin{problem}\label{pr:P1-average}
Bounding the norm of each section \(s^al^{D-a}\) by \(1\) for unit vectors \(s, l\), and averaging over \(s, l\) in the unit spheres, gives \(\deg E_D \ge -\frac14 D(D+1)\).  For \(D = 10\) this is \(-27.5\), against \(\deg E_{10} = 32.1969\ldots\)
\end{problem}

So at an orthonormal pair nothing is lost, and the main term \(\frac14D^2\) of Theorem~\ref{cv-thm:P1} comes from the small norms of the columns, about \(\binom{D}{a}^{-1/2}\).  Averaging over \(s\) and \(l\) loses \(\frac12 D^2\), because it spreads the pair over pairs whose sections are far from orthogonal.  For a curve \(C \subset \Pb^2\) no orthogonal family of this kind is known, and the argument described in \cite[proof dependency map, item HSAt]{Adr} uses Chow forms and a determinant estimate.  The argument needs a lower bound of this kind, with explicit constants, on varieties of growing dimension, where Section~\ref{cv-sec:uniform} recorded that such bounds lose a factor exponential in the dimension.

\section*{Hints}
\paragraph{Section~\ref{sec:mahler}.}
\begin{description}
\item[\ref{pr:jensen}] Jensen's formula for \(x - \alpha\).  For \(\beta_i = 0\) the factor is a constant.
\item[\ref{pr:lex}] Induction on \(N\), using Problem~\ref{pr:jensen} in the last variable for fixed values of the others.
\item[\ref{pr:beta}] Write \(z = g/\lvert g \rvert_2\) with \(g\) a standard complex Gaussian vector; \(\lvert g_k \rvert^2\) are independent exponential variables.
\item[\ref{pr:sphere-torus}] The measure on the sphere is invariant under the torus.  For fixed moduli \(r_k = \lvert u_k \rvert\) apply Problem~\ref{pr:lex} to \(G(r_1x_1, \dots, r_Nx_N)\), then average over the moduli with Problem~\ref{pr:beta}.
\item[\ref{pr:sphere-linear}] Unitary invariance moves \(\bar p\) to a multiple of the first basis vector.
\end{description}
\paragraph{Section~\ref{sec:points}.}
\begin{description}
\item[\ref{pr:point}] Problems~\ref{pr:sphere-linear} and~\ref{pr:jensen}; \(m(F) = m(F(x,1))\), since \(F\) is homogeneous.
\item[\ref{pr:quotient0}] Gauss's lemma.  Evaluation at the \(\omega_i\), divided by \(\lvert\omega_i\rvert_2^D\), is an isometry of \(\Lambda_D \otimes \C\) onto \(\C^\delta\).
\item[\ref{pr:index0}] Order the monomials by the power of \(x\) and add the columns \(F x^k y^{D-\delta-k}\).
\item[\ref{pr:vdm}] Factor \(\omega_{i1}^{D-\delta+1}\) out of row \(i\).  Write \(\omega_i = t_i(\alpha_i, 1)\) with \(\prod t_i = a\).
\item[\ref{pr:HS0}] Theorem~\ref{thm:points}.
\end{description}
\paragraph{Section~\ref{sec:chow}.}
\begin{description}
\item[\ref{pr:chow-curve}] For fixed \(s\), \(l \mapsto s \times l\) maps the lines onto the points of the line \(s = 0\).
\item[\ref{pr:chow-line}] Fix \(s\) and apply Problem~\ref{pr:sphere-linear} in \(l\); then average \(\log(\lvert s_0 \rvert^2 + \lvert s_1 \rvert^2)\) with Problem~\ref{pr:beta}.
\item[\ref{pr:line-deg}] Substitute \(x_2 = Nx_0\) for the basis.  For the norm, put \(u := \sqrt{1 + N^2}\, x_0\) and \(v := x_1\).  On \(C_N\), \(\lvert z \rvert_2^2 = \lvert u \rvert^2 + \lvert v \rvert^2\), so \((u : v)\) identifies \(C_N(\C)\), with its metric, with \(\Pb^1(\C)\).
\item[\ref{pr:line-chow}] Problem~\ref{pr:chow-curve}, then a matrix in \(\mathrm{SU}(3)\) taking \(p/\lvert p \rvert_2\) to \((0, 0, 1)\), applied to \(s\) and \(l\); then Problem~\ref{pr:chow-line} and~\eqref{eq:philippon}.
\end{description}
\paragraph{Section~\ref{sec:check}.}
\begin{description}
\item[\ref{pr:P1-det}] The matrix is \(\mathrm{Sym}^D\) of a \(2 \times 2\) matrix; check diagonal matrices first.
\item[\ref{pr:P1-exact}] The Fubini--Study norm is invariant under \(\mathrm{U}(2)\), which moves \((s, l)\) to \((y, x)\).
\item[\ref{pr:P1-average}] Hadamard's inequality in \(E_D \otimes \C\), with \(\lvert s(z) \rvert \le \lvert z \rvert_2\); then Problems~\ref{pr:P1-det}, \ref{pr:chow-line} and~\ref{cv-pr:bergman}.
\end{description}

\begin{thebibliography}{GKZ}
\bibitem[Adr]{Adr} J.~Adriazola, \emph{The Alp\"oge proof of the Riemann hypothesis}, manuscript, revised 25 November 2026, \href{https://github.com/jimjamzola/RH-Alpoge-Reconstruction/tree/9eefc08b0819728d4dd2ef0568b49a61251af111}{commit \texttt{9eefc08}}.
\bibitem[BG]{BG} E.~Bombieri and W.~Gubler, \emph{Heights in Diophantine Geometry}, New Math.\ Monogr.\ 4, Cambridge Univ.\ Press, 2006.
\bibitem[BGS]{BGS} J.-B.~Bost, H.~Gillet and C.~Soul\'e, \emph{Heights of projective varieties and positive Green forms}, J.\ Amer.\ Math.\ Soc.\ \textbf{7} (1994), 903--1027.
\bibitem[GKZ]{GKZ} I.~M.~Gelfand, M.~M.~Kapranov and A.~V.~Zelevinsky, \emph{Discriminants, Resultants, and Multidimensional Determinants}, Birkh\"auser, 1994.
\bibitem[KPS]{KPS} T.~Krick, L.~M.~Pardo and M.~Sombra, \emph{Sharp estimates for the arithmetic Nullstellensatz}, Duke Math.\ J.\ \textbf{109} (2001), 521--598.
\bibitem[Le]{Le} P.~Lelong, \emph{Mesure de Mahler et calcul de constantes universelles pour les polynomes de \(N\) variables}, Math.\ Ann.\ \textbf{299} (1994), 673--695.
\bibitem[NP]{NP} Yu.~V.~Nesterenko and P.~Philippon (eds.), \emph{Introduction to Algebraic Independence Theory}, Lecture Notes in Math.\ 1752, Springer, 2001.
\bibitem[Ph1]{Ph1} P.~Philippon, \emph{Sur des hauteurs alternatives.\ I}, Math.\ Ann.\ \textbf{289} (1991), 255--283.
\bibitem[Ph3]{Ph3} P.~Philippon, \emph{Sur des hauteurs alternatives.\ III}, J.\ Math.\ Pures Appl.\ (9) \textbf{74} (1995), 345--365.
\end{thebibliography}

\end{document}
